\section{Phần D --- Ứng dụng: trợ lý thủ tục hành chính}
\label{sec:phan-d}

Phần~D là bài toán Việt Nam, trọng số 30 điểm, gồm năm hạng mục D.1--D.5
\cite[tr.~19]{de}. Mục này trình bày \emph{trạng thái thực tế}: những gì đã làm có
bằng chứng trong kho mã, và những ngưỡng của đề chưa đạt được. Nhóm không trình
bày kế hoạch như kết quả.

\textbf{Tóm tắt trạng thái.} D.1 \motphan; D.2 \motphan\ (400 QA ứng viên, 24 đã
duyệt); D.3 \motphan\ (pilot 24 QA đã chạy trên Kaggle, chưa chạy tập chính);
D.4 \chualam; D.5 \chualam.

%--------------------------------------------------------------------
\subsection{D.1 --- Thu thập mô tả thủ tục công khai}
\label{sec:d1}

\textbf{Yêu cầu.} Thu thập trình tự, hồ sơ, thời hạn, lệ phí và căn cứ pháp lý từ
Cổng Dịch vụ công quốc gia và cổng địa phương; lưu truy xuất và phiên bản.

\subsubsection{Năm họ thủ tục đã rà \motphan}

\begin{center}
\small
\begin{tabular}{p{0.40\textwidth}lp{0.28\textwidth}}
\toprule
\textbf{Họ thủ tục} & \textbf{Mã} & \textbf{Trạng thái căn cứ}\\
\midrule
Cấp bản sao trích lục hộ tịch / bản sao giấy khai sinh & \texttt{2.000635} &
Đã xác minh QĐ 1818 và QĐ 2163; \emph{thời hạn chưa thống nhất}\\
Chứng thực bản sao từ bản chính & \texttt{2.000815} &
Đã xác minh Điều 20 NĐ 23/2015 hợp nhất 753/VBHN-BTP\\
Chứng thực chữ ký & \texttt{2.000884} &
Đã xác minh phụ lục QĐ 858; chưa chốt quy trình tại Hà Nội\\
Thay đổi, cải chính, bổ sung hộ tịch, xác định lại dân tộc & \texttt{1.004859} &
Đã xác minh QĐ 163 và biểu lệ phí NQ 98/2026/NQ-HĐND\\
Xác nhận thông tin hộ tịch & \texttt{2.002516} &
Đã xác minh QĐ 1852; chuyển dẫn chứng từ video sang văn bản\\
\bottomrule
\end{tabular}
\end{center}

Mỗi bản ghi lưu URL, thời điểm truy xuất và SHA-256 của bản tải. Nhóm cũng kiểm
chữ ký CMS về toàn vẹn nội dung trên bốn PDF Hà Nội --- đây là kiểm toàn vẹn tệp,
\emph{không} phải thẩm định chuỗi tin cậy.

\subsubsection{Hai khác biệt phải xử lý trước khi huấn luyện}

Quá trình rà phát hiện hai mâu thuẫn mà nếu bỏ qua sẽ sinh dữ liệu sai:

\begin{enumerate}[leftmargin=2em]
\item \textbf{Thời hạn cấp trích lục không thống nhất giữa các văn bản.} Phụ lục
QĐ 1818/QĐ-TTPVHCC (hiệu lực 01/01/2026) ghi \emph{dự kiến 3 ngày làm việc}, tối đa
10 ngày nếu phải xác minh. QĐ 2163/QĐ-BTP (30/09/2026) ghi \emph{trong ngày tiếp
nhận, hoặc ngày làm việc kế tiếp nếu sau 15 giờ}. Nhóm \emph{chưa} tìm được văn bản
triển khai 2163 tại Hà Nội, cũng chưa tìm được văn bản xác nhận 1818 đã bị thay
thế. Xử lý: không dùng ``3/10 ngày'' làm đáp án chung hiện hành; chỉ soạn câu hỏi
gắn với \emph{từng văn bản và phạm vi của nó}.
\item \textbf{Lệ phí không phải con số hiển thị trên cổng.} Thông tư
226/2016/TT-BTC là căn cứ phí \emph{chứng thực}, không phải biểu lệ phí hộ tịch Hà
Nội; biểu hiện hành của Hà Nội nằm ở Nghị quyết 98/2026/NQ-HĐND. Trộn hai khái niệm
này tạo ra đáp án sai một cách khó phát hiện.
\end{enumerate}

Hai phát hiện này minh họa vì sao hạng mục D.1 không thể rút gọn thành ``cào trang
thủ tục'': nội dung hiển thị trên cổng không đồng nhất với văn bản có hiệu lực.

\subsubsection{Căn cứ quyền sử dụng}

Nhóm phân loại nội dung trước khi đưa vào dữ liệu:

\begin{itemize}[leftmargin=2em]
\item \textbf{Chính văn quy phạm và văn bản hành chính do cơ quan ban hành:} theo
khoản 2 Điều 15 Luật Sở hữu trí tuệ (VBHN 160/2026) và khoản 2 Điều 8 VBHN
6154/2026, nhóm văn bản này \emph{không} thuộc phạm vi bảo hộ quyền tác giả. Có
căn cứ để dùng phần chính văn công khai đã xác minh làm nguồn soạn QA, kèm số hiệu,
điều/mục, phiên bản và URL.
\item \textbf{Trang chi tiết của cổng, bài tin, bài hướng dẫn, video, bản lưu
pilot:} dùng để \emph{tìm} văn bản và đối chiếu cách hiển thị. Điều khoản Cổng quốc
gia yêu cầu dùng lại đúng pháp luật; footer cổng Hà Nội ghi
\emph{All Rights Reserved} và yêu cầu ghi nguồn. Nhóm \textbf{chưa} chốt căn cứ để
lấy nguyên văn toàn trang hoặc phân phối lại chúng, nên khi cần cùng một quy định
thì trích từ văn bản chính thức, không chuyển snapshot web sang tập train.
\end{itemize}

Về truy cập: \texttt{robots.txt} của cổng Hà Nội cho \texttt{Allow: /} với
\texttt{Crawl-delay: 5}, nên nhóm truy xuất tuần tự cách ít nhất 5 giây và không
crawl toàn cổng. Với các host trả HTTP 404 cho \texttt{robots.txt}, nhóm ghi rõ
rằng \emph{không có tệp robots không đồng nghĩa có giấy phép nội dung}.

\textbf{Còn thiếu.} Chưa phát hành bộ nguồn huấn luyện; còn phải kiểm các kênh và
phiên bản triển khai thực tế tại Hà Nội.

%--------------------------------------------------------------------
\subsection{D.2 --- Bộ dữ liệu chỉ dẫn có trích dẫn}
\label{sec:d2}

\textbf{Yêu cầu của đề.} Tạo $\ge 2.000$ QA có trích dẫn bằng phương pháp bán tự
động; kiểm thủ công $\ge20\%$; tập test là câu hỏi của \emph{người dùng thực} thu
qua khảo sát, tách hoàn toàn khỏi train.

\subsubsection{Đã làm \motphan}

\begin{itemize}[leftmargin=2em]
\item \textbf{Quy trình sinh có cổng kiểm.} Nhóm viết recipe soạn QA từ chính văn,
kiểm SHA-256 và từng đoạn trích, CLI chọn mẫu phân tầng, và một cổng xuất SFT từ
chối phát hành khi dữ liệu chưa đạt điều kiện.
\item \textbf{Công cụ chặn rò rỉ test.} Kiểm tách \texttt{USER\_TEST}, chặn đường
dẫn thiếu, trùng họ thủ tục và trùng câu.
\item \textbf{Lô pilot chính văn v2 đã duyệt và phát hành nội bộ:} 24 QA, nhóm
kiểm thủ công \textbf{24/24}, có biên bản chung của lô. Cổng D.2 trả
\texttt{pilot\_sft\_ready}; đầu vào D.3 kiểm \texttt{release-valid} với 18 train
và 6 validation. Đây là lô duy nhất đã phát hành.
\item \textbf{Lô pilot v1 trước đó:} 20 QA, 21 đoạn trích, 9 nguồn; đã rà riêng.
Validator ngoại tuyến xác nhận 20/20 qua kiểm cấu trúc và 21/21 đoạn trích khớp
bản lưu --- nhưng kiểm tự động không được tính vào quota kiểm thủ công.
\item \textbf{Tập chính đang mở rộng theo từng lô:} b01 thêm 57 QA train, b02
thêm 54, b03 thêm 265. Tích lũy \textbf{394 QA train và 6 validation ứng viên},
\emph{chưa phát hành}. Các lô mới gồm cả câu cần từ chối hoặc nói rõ giới hạn
nguồn.
\end{itemize}

\subsubsection{Khoảng cách tới ngưỡng của đề}

\begin{center}
\begin{tabular}{lrrl}
\toprule
\textbf{Ngưỡng} & \textbf{Yêu cầu} & \textbf{Hiện có} & \textbf{Trạng thái}\\
\midrule
Số QA có trích dẫn & $\ge2.000$ & 400 ứng viên & \chualam\\
Kiểm thủ công & $\ge20\%$ ($\ge400$) & 24 đã duyệt & \chualam\\
Tập test từ khảo sát người thật & có & chưa có & \chualam\\
Tách train/validation/test & có & có cho lô pilot & \motphan\\
Datasheet & có & đã soạn khung & \motphan\\
\bottomrule
\end{tabular}
\end{center}

\textbf{Vì sao tập test không thể sinh bằng máy.} Đề yêu cầu test là câu hỏi của
người dùng thực. Mục đích là đo khả năng tổng quát hóa sang \emph{cách người dân
thật diễn đạt} --- thường thiếu thuật ngữ, thiếu bối cảnh, hoặc mô tả tình huống cá
nhân thay vì hỏi tên thủ tục. Dùng AI sinh câu hỏi thay khảo sát sẽ phá chính phép
đo đó, vì câu hỏi sinh ra mang phân phối của mô hình, không phải của người dùng.
Quy tắc nội bộ của nhóm cấm việc này.

Thu dữ liệu người cần văn bản đồng ý và ẩn danh trước khi lưu
(Mục~\ref{sec:dao-duc}), nên hạng mục này có một \emph{sàn thời gian} không nén
được bằng công sức.

%--------------------------------------------------------------------
\subsection{D.3 --- Tinh chỉnh có giám sát bằng QLoRA}
\label{sec:d3}

\textbf{Yêu cầu.} Chọn LM mã nguồn mở hỗ trợ tiếng Việt, 0,5--3B tham số, giấy phép
cho nghiên cứu; SFT bằng LoRA hoặc QLoRA; tối ưu $r$, $\alpha$, learning rate,
epochs \emph{chỉ trên validation}.

\subsubsection{Đã chuẩn bị \motphan}

\begin{itemize}[leftmargin=2em]
\item \textbf{Mô hình:} Qwen2.5-1.5B-Instruct, Apache-2.0, revision được khóa
trong \texttt{scripts/}. Nằm trong khoảng 0,5--3B và hỗ trợ tiếng Việt. Giấy phép
và revision đã đối chiếu với model card \cite{qwen}.
\item \textbf{Mã:} bộ nạp dữ liệu có che prompt, mã QLoRA, notebook Kaggle, script
đóng gói tối thiểu, và kế hoạch so cấu hình $r$, $\alpha$, LR, epoch trên
validation.
\item \textbf{Cổng an toàn đã kiểm:} ZIP chỉ chứa mã thì không train; sai hash thì
dừng; CLI từ chối manifest \texttt{candidate\_only} của lô 24 QA. Nghĩa là quy
trình \emph{không thể} vô tình huấn luyện trên dữ liệu chưa được duyệt.
\item \textbf{Kiểm phần mềm trên CPU:} 7 test bộ nạp và mặt nạ đều qua. Với một câu
minh họa dùng tokenizer Qwen đúng revision: 71 token prompt bị che, giải mã 19
token đích khớp nguyên đáp án có dấu và trích dẫn, đúng một EOS. Đã chạy forward
trên một Qwen thu nhỏ với trọng số khởi tạo và adapter Q/V.
\item \textbf{Kiểm resume:} với mô hình thu nhỏ, checkpoint khôi phục đúng trọng số
adapter, trạng thái AdamW, RNG CPU và thứ tự trộn.
\end{itemize}

\subsubsection{Kết quả pilot đã chạy trên Kaggle \dalam}

Nhóm đã chạy một lượt SFT pilot thật trên lô 24 QA đã phát hành. Số liệu dưới đây
đối chiếu từ \texttt{metrics.jsonl} và \texttt{summary.json} do nhóm gửi.

\begin{center}
\begin{tabular}{lrr}
\toprule
\textbf{Epoch} & \textbf{Validation NLL} & \textbf{Validation PPL}\\
\midrule
1 & 1,509630 & 4,525054\\
2 & \textbf{1,242314} & \textbf{3,463617}\\
\bottomrule
\end{tabular}
\end{center}

Seed 41 hoàn thành $2/2$ epoch, không dừng sớm; mỗi epoch giám sát 1.235 token
train và 388 token validation. Epoch 2 được chọn theo validation NLL. PPL khớp
$e^{\mathrm{NLL}}$ theo \eqref{eq:b6-ppl}.

\textbf{Cấu hình đã chạy:} QLoRA NF4 double quantization, FP16, adapter ở Q/V,
$r=8$, $\alpha=16$, dropout $0$, learning rate $10^{-4}$, batch $1$, độ dài tối đa
$512$. Môi trường: Tesla T4, Python 3.13.16, torch 2.11.0+cu130, CUDA 13.0.

\textbf{Bộ nhớ đã đo:} peak allocated $2{,}402$ GiB, peak reserved $2{,}725$ GiB
theo PyTorch. Đây là số của tiến trình, không phải tổng VRAM toàn thiết bị. Con số
này nay \emph{thay thế} phần ước lượng ở Mục~\ref{sec:c3} cho đúng cấu hình trên.

\textbf{Tái lập đã kiểm:} bốn SHA-256 (manifest release-01, hai tệp Python trong
ZIP mã, chat template tokenizer) đều khớp. Resume từ \texttt{epoch-0001} tái lập
đúng epoch 2: cùng 1.235 token train (train NLL $1{,}750317$), cùng 388 token
validation, NLL $1{,}242314$ và PPL $3{,}463617$; chỉ hai trường peak VRAM lệch.
Resume \emph{không} được tính là một seed độc lập.

%--------------------------------------------------------------------
\subsubsection{Hỏi thử adapter pilot: hai lỗi quan sát được}

Nhóm chạy một notebook hỏi thử riêng, suy luận từ \texttt{best\_adapter/} và
\emph{không} huấn luyện lại. Mô hình sinh được hai đáp án, nhưng \textbf{cả hai
đều không đạt} mục tiêu đặt ra ở Mục~\ref{sec:gioi-thieu}:

\begin{enumerate}[leftmargin=2em]
\item \textbf{Sai thời hạn và tự nhận là mới nhất.} Với câu hỏi về xác nhận
thông tin hộ tịch, mô hình trả ``15/30 ngày làm việc'' và tự nhận đây là quy định
mới nhất. Phụ lục 1852/QĐ-TTPVHCC, B.4 tr.~3 ghi \emph{dự kiến 03 ngày}, và không
quá 10 ngày làm việc nếu cần kiểm tra, xác minh.
\item \textbf{Không từ chối câu ngoài phạm vi.} Với câu hỏi về thị thực du lịch
Nhật Bản --- hoàn toàn ngoài phạm vi hộ tịch và chứng thực tại Hà Nội --- mô hình
\emph{hướng dẫn} thay vì nêu rõ là ngoài phạm vi.
\end{enumerate}

Thêm nữa, cả hai đáp án chỉ ghi \texttt{[1]} mà không kèm nguồn truy xuất được.

\textbf{Vì sao hai lỗi này đáng ghi vào báo cáo.} Lỗi thứ hai chính là hiện tượng
mà Mục~\ref{sec:d5} đặt giả thuyết: SFT dạy mô hình \emph{luôn} trả lời, nên
``từ chối đúng cách'' là thứ SFT khó dạy. Quan sát này \emph{ủng hộ} giả thuyết
nhưng không chứng minh nó, và nhất định không phải điểm D.5.

\textbf{Đây là quan sát định tính, không phải phép đo.} Hai câu, một adapter huấn
luyện trên 24 QA, không có người đánh giá mù. Nhóm cũng \emph{không} chuyển hai QA
validation cùng dẫn đoạn B.4 sang train để làm đẹp kết quả demo --- làm vậy là rò
rỉ validation.

\subsubsection{Hạng mục D.3 còn lại \chualam}

\begin{itemize}[leftmargin=2em]
\item \textbf{Chưa chọn cấu hình.} Đề yêu cầu tối ưu $r$, $\alpha$, learning rate
và epochs trên validation. Hiện chỉ có \emph{một} cấu hình đã chạy ($r=8$,
$\alpha=16$, LR $10^{-4}$), nên chưa có việc chọn nào diễn ra.
\item \textbf{Chưa chạy trên tập chính.} Lượt pilot dùng 24 QA; 394 QA train của
tập chính chưa phát hành nên chưa huấn luyện.
\item \textbf{Chưa đủ 5 lần chạy độc lập.} Chỉ có seed 41. Yêu cầu tái lập ở
Mục~\ref{sec:tai-lap} cần ít nhất 5 seed; resume không tính là seed thứ hai.
\end{itemize}

\subsection{D.4 --- Dữ liệu ưu tiên và huấn luyện DPO}
\label{sec:d4}

\textbf{Yêu cầu.} Thu tối thiểu \textbf{500 cặp} so sánh: hai đáp án SFT cho cùng
một câu hỏi; \textbf{ít nhất hai người đánh giá độc lập} theo đúng/đầy đủ/trích
dẫn; báo cáo mức đồng thuận; huấn luyện DPO.

\textbf{Trạng thái: \chualam.}

\subsubsection{Quy trình đã thiết kế}

Nhóm đã viết hướng dẫn chấm và rubric, và chốt quy trình:

\begin{enumerate}[leftmargin=2em]
\item Với mỗi câu hỏi \emph{ngoài tập test}, sinh hai đáp án từ checkpoint SFT
\emph{thật}. Khóa checkpoint, template và cấu hình sinh trước khi chấm.
\item Hà Mạnh Thế và Trần Duy Hưng chấm \emph{độc lập} theo ba tiêu chí: đúng nội
dung, đủ ý, trích dẫn đúng. \textbf{Không chia đôi 500 cặp cho hai người} --- mỗi
người chấm toàn bộ.
\item Không ai được xem nhãn của người kia trước khi khóa.
\item Khóa nhãn, \emph{đo mức đồng thuận}, rồi Cao Văn Hà mới phân xử. Giữ nguyên
nhãn gốc và lý do loại cặp; không ghi đè nhãn raw.
\item Dùng các cặp ưu tiên để huấn luyện DPO, giữ mô hình SFT tham chiếu cố định,
và chọn $\beta$ \emph{trên validation DPO}.
\end{enumerate}

Thiết kế này trực tiếp phản ánh Mục~\ref{sec:b5}: $\pref$ phải là một mô hình xác
định và cố định, vì \eqref{eq:b5-gibbs} định nghĩa nghiệm tối ưu \emph{tương đối}
với nó; hai đáp án phải thuộc cùng một prompt để hạng $\beta\log Z(x)$ triệt tiêu
ở \eqref{eq:b5-invert}.

\subsubsection{Hai lý do hạng mục này bị chặn}

\begin{enumerate}[leftmargin=2em]
\item \textbf{Chặn kỹ thuật --- đã gỡ một phần.} Trước đây không có checkpoint SFT
nào, nên không sinh được hai đáp án để so sánh. Nay đã có adapter pilot
(\texttt{best\_adapter/}, seed 41, epoch 2), nên \emph{về mặt kỹ thuật} nhóm đã
sinh được cặp đáp án. Giới hạn còn lại là chất lượng: adapter này huấn luyện trên
24 QA và hỏi thử cho thấy nó sai thời hạn và không từ chối câu ngoài phạm vi, nên
cặp sinh từ nó chưa đại diện cho mô hình cuối. Nhóm có thể dùng nó để \emph{hiệu
chuẩn rubric và tập chấm thử}, nhưng 500 cặp chính thức nên lấy từ checkpoint huấn
luyện trên tập chính.
\item \textbf{Chặn về khối lượng người:} yêu cầu ``hai người đánh giá độc lập'' có
nghĩa là \emph{hai con người}, mỗi người chấm toàn bộ số cặp. Phân tích tải nội bộ
của nhóm (dùng giả định 2 phút/cặp, không phải số đo) cho thấy 700 cặp thô cần
khoảng 23 giờ cho mỗi người chấm. Đây là một \emph{sàn thời gian}, không phải vấn
đề hiệu suất.
\end{enumerate}

\textbf{Điều nhóm không làm.} Không sinh nhãn bằng AI để thay người chấm. Nếu làm
vậy, con số ``mức đồng thuận'' sẽ không đo gì cả, và toàn bộ D.4 cùng D.5 (dựa lên
D.4) mất giá trị. Quy tắc nội bộ ghi rõ: không bịa điểm người chấm hay mức đồng
thuận.

\subsubsection{Hệ quả cho cầu nối A$\leftrightarrow$D (G5)}

Đề yêu cầu ít nhất một cài đặt Phần~A được dùng hoặc so sánh trực tiếp trong
Phần~D. Thiết kế của nhóm là đối chiếu \texttt{src/dpo.py} (A.4) với cách tính
chuẩn trên \emph{một minibatch dữ liệu ưu tiên thật} của D.4, so cả mất mát và
gradient. Phép đối chiếu này \textbf{chưa} thực hiện được vì chưa có dữ liệu đó.
Đây là hạng mục còn thiếu rõ ràng nhất của báo cáo.

%--------------------------------------------------------------------
\subsection{D.5 --- Đánh giá ba mốc mô hình}
\label{sec:d5}

\textbf{Yêu cầu.} So sánh base, SFT, SFT+DPO về độ chính xác đối chiếu nguồn,
hallucination, và từ chối khi ngoài phạm vi; blind human pairwise evaluation với
win rate kèm khoảng tin cậy; nêu giới hạn khi kiến thức thay đổi nhanh và hướng
RAG.

\textbf{Trạng thái: \chualam.}

\subsubsection{Giao thức dự kiến}

\begin{itemize}[leftmargin=2em]
\item Ba mốc model đánh giá trên \emph{cùng} tập test (tập khảo sát người thật của
D.2).
\item Ba nhóm chỉ số có định nghĩa trước: (i) đáp án nêu đúng điều kiện áp dụng và
dẫn đúng nguồn; (ii) tỷ lệ phát biểu không có căn cứ trong nguồn đã dẫn;
(iii) tỷ lệ từ chối đúng khi câu hỏi ngoài phạm vi hoặc chưa đủ căn cứ.
\item \textbf{Người đánh giá không biết đáp án thuộc model nào} (blind), thứ tự
trình bày được xáo.
\item Win rate báo cáo kèm khoảng tin cậy, ghi rõ đơn vị lấy mẫu.
\end{itemize}

\subsubsection{Vì sao chỉ số thứ ba quan trọng nhất}

Chỉ số (iii) đo đúng điều kiện thứ ba trong đặc tả sản phẩm
(Mục~\ref{sec:gioi-thieu}). Một mô hình đạt điểm cao ở (i) nhưng thấp ở (iii) là mô
hình nguy hiểm: nó trả lời tự tin cả khi không có căn cứ. Giả thuyết của nhóm là
DPO giúp nhiều nhất ở (iii), vì ``từ chối đúng cách'' là thứ khó dạy bằng SFT ---
dữ liệu SFT gồm các cặp câu hỏi--đáp án, nên nó dạy mô hình \emph{luôn} trả lời.

Đây là \emph{giả thuyết}, không phải kết quả. Nhóm không hứa DPO tốt hơn SFT và sẽ
báo cả kết quả âm nếu có.

\subsubsection{Phụ thuộc và giới hạn}

D.5 cần cả ba mốc model (phụ thuộc D.3 và D.4), tập test khảo sát thật (phụ thuộc
D.2), và người đánh giá mù --- là người thứ ba ngoài nhóm hoặc thành viên không
tham gia tạo nhãn. Không có đủ các điều kiện này thì bảng ba hệ không thể lập.

Về RAG: Đề tài 11 trong danh mục của đề là hướng truy xuất tăng cường. Nhóm chỉ
\emph{thảo luận} đây là hướng mở rộng hợp lý cho vấn đề kiến thức thay đổi nhanh
(xem hai mâu thuẫn ở Mục~\ref{sec:d1}), và \textbf{không} trình bày RAG như một
hạng mục đã triển khai.

%--------------------------------------------------------------------
\subsection{Tổng hợp trạng thái Phần D}

\begin{center}
\begin{tabular}{llp{0.46\textwidth}}
\toprule
\textbf{Mã} & \textbf{Trạng thái} & \textbf{Khoảng cách còn lại}\\
\midrule
D.1 & \motphan & Chưa phát hành bộ nguồn; còn kiểm kênh/phiên bản tại Hà Nội\\
D.2 & \motphan & 400/2.000 QA ứng viên; 24/400 kiểm thủ công; chưa có test khảo sát thật\\
D.3 & \motphan & Pilot 24 QA đã chạy (PPL validation 3,46; peak 2,402 GiB); chưa chọn cấu hình, chưa chạy tập chính, chưa đủ 5 seed\\
D.4 & \chualam & Đã có adapter pilot để sinh cặp; cần hai người chấm toàn bộ $\ge500$ cặp\\
D.5 & \chualam & Bị chặn bởi D.4 và tập test khảo sát; cần người đánh giá mù\\
\midrule
G5 & \chualam & Chưa đối chiếu A.4 trên minibatch ưu tiên thật của D\\
\bottomrule
\end{tabular}
\end{center}

\textbf{Đánh giá trung thực.} Phần~D chưa đạt, nhưng đường ống đã chạy hết một
vòng ở quy mô pilot: dữ liệu có dẫn nguồn $\to$ cổng phát hành $\to$ SFT trên GPU
$\to$ adapter sinh được đáp án. Điều này chứng minh thiết kế hoạt động.

Phần còn lại bị chặn bởi hai thứ khác nhau. Thứ nhất là \emph{khối lượng}: 394 QA
train cần người duyệt, và tập chính cần huấn luyện lại với nhiều cấu hình và $\ge5$
seed. Thứ hai, và không nén được bằng công sức của mã, là \emph{người thật}: người
khảo sát có văn bản đồng ý cho tập test (D.2), hai người chấm độc lập cho
$\ge500$ cặp (D.4), và người đánh giá mù (D.5). Nhóm chọn báo cáo đúng trạng thái
này thay vì điền số liệu không có nguồn.
